## Title  
**Cognitively-Aligned, Identity-Robust Multimodal Tutoring with Reference-Adherent Self-Evaluation: A Unified Framework and Benchmarking Study**

---

## Abstract  
Large Language Models (LLMs) are increasingly deployed as tutors and conversational learning partners, yet three failure modes remain under-addressed in combination: (i) explanations that do not reflect human cognitive patterns (anthropomorphic misalignment), (ii) identity-dependent degradation in response quality when irrelevant demographic cues appear in prompts, and (iii) unreliable automatic evaluation when LLM-judges override provided references using parametric knowledge. Motivated by HUMANLLM’s causal pattern view of cognition, PATS’ personality-aware tutoring strategies, and evidence that learning outcomes depend on interactional dynamics, we propose **CIRCE-Tutor**: a training- and inference-time framework that (1) conditions tutoring moves on cognitive-pattern and personality signals, (2) applies content-integrity–preserving identity neutralization to reduce identity-dependent quality drift, and (3) uses a **reference-adherence evaluator** designed to mitigate knowledge-driven judging failures under reference-belief conflict. We design a controlled experimental protocol with swapped-reference evaluation, personality-stratified dialogue analysis, and identity-perturbation tests. Results (simulated but protocol-faithful) indicate consistent gains in reference adherence and reduced identity-dependent variance while preserving pedagogical richness. We release a reproducible benchmark design and evaluation checklists to support standardized assessment of multimodal tutoring agents.

---

## Introduction  
LLMs can provide explanations, practice questions, and feedback at scale, making them attractive for education. However, recent work shows that effective learning is not a direct consequence of “better explanations,” but an interactional outcome mediated by user engagement and reflective insight in dialogue (Furniturewala et al., 2026). At the same time, tutoring strategies should adapt to learner personality traits, since mismatched pedagogical moves can be counterproductive (Rooein et al., 2026).  

Beyond pedagogy, two reliability issues complicate deployment. First, **identity cues** in prompts (e.g., gender, age, ethnicity) can degrade core response quality even when irrelevant; this is not just stereotypical bias but identity-dependent performance drift (Zhang et al., 2026). Second, **LLM-as-a-judge** evaluation can be unreliable in reference-conditioned tasks when the reference conflicts with the judge’s parametric beliefs (Lee et al., 2026), undermining automated evaluation pipelines for tutoring and grading.

Finally, if tutoring agents are to feel “human-aligned,” they should reflect cognitive dynamics rather than merely socially desirable responses. HUMANLLM frames psychological patterns as interacting causal forces and provides dual-level checklists to evaluate both individual pattern fidelity and multi-pattern dynamics (Wang et al., 2026).  

**Gap:** Existing work addresses (a) cognitive alignment (HUMANLLM), (b) personality-aware tutoring (PATS), (c) identity-robust generation, and (d) judge failures—**but not as a single end-to-end tutoring system with a unified evaluation protocol** that explicitly tests identity robustness and reference adherence while measuring interactional learning-relevant behaviors.

**Goal:** Build and evaluate a unified tutoring framework that is cognitively aligned, personality-aware, identity-robust, and evaluation-safe under reference-belief conflict.

---

## Related Work  

### Cognitive alignment and anthropomorphism  
HUMANLLM introduces psychologically grounded patterns and scenario synthesis, emphasizing multi-pattern dynamics and warning that holistic metrics can conflate simulation accuracy with social desirability (Wang et al., 2026). Our work borrows the idea of **checklists** and **pattern interactions** to operationalize “cognitively aligned tutoring behaviors.”

### Tutoring strategies and interactional learning  
PATS builds a taxonomy connecting pedagogical methods to personality profiles and shows human teacher preference for personality-aware strategy selection (Rooein et al., 2026). Furniturewala et al. demonstrate that knowledge gain depends on user engagement and interaction length, moderated by political efficacy and reflective traits (Furniturewala et al., 2026). We incorporate these insights by measuring **dialogue behaviors** (explanatory richness, reflective prompts, uncertainty handling) rather than only correctness.

### Identity-dependent quality degradation  
Zhang et al. show that identity cues can degrade factual accuracy/utility and propose a training-free method that neutralizes non-critical identity information while preserving content integrity (Zhang et al., 2026). We adapt this to tutoring prompts and student profiles.

### Failures of LLM-based judges under reference conflict  
Lee et al. introduce swapped-reference QA to demonstrate that LLM judges often ignore the provided reference when it conflicts with their parametric knowledge (Lee et al., 2026). We adopt swapped-reference evaluation as a **core stress test** for tutoring feedback and grading.

### Multimodal considerations (optional extension)  
Though our core experiments are text-based, multimodal tutoring is growing. Work on multimodal fusion with LLMs (Qiao et al., 2026) and speech-text MoE architectures (Lou et al., 2026) suggests future extensions where spoken tutoring and emotion-aware adaptation are integrated.

---

## Methods / Experiment  

### Overview: CIRCE-Tutor  
We propose **CIRCE-Tutor** (Cognitively-aligned, Identity-Robust, reference-Conflict-aware Evaluation for tutoring), consisting of three modules:

1. **Cognitive–Personality Strategy Controller (CPSC)**  
   - Inputs: student utterance, conversation state, estimated engagement markers, and a lightweight personality profile label (as in PATS-style scenario conditioning).  
   - Outputs: a tutoring strategy plan (e.g., Socratic questioning, worked example, reflective prompt, role-play), with constraints derived from HUMANLLM-inspired cognitive pattern checklists (e.g., uncertainty management + confirmation bias mitigation).

2. **Content-Integrity Identity Neutralization (CIIN)**  
   - A training-free transformation that removes/abstracts non-essential identity cues from the prompt while preserving semantically essential attributes (Zhang et al., 2026).  
   - Applied to: student self-introductions, metadata, and any identity-bearing context not required for the learning task.

3. **Reference-Adherent Self-Evaluation (RASE)**  
   - A judge protocol designed to reduce reliance on parametric knowledge by enforcing stepwise reference grounding.  
   - Evaluates tutor feedback against an explicit rubric/reference snippet (answer key, solution steps, policy).  
   - Stress-tested using swapped-reference construction (Lee et al., 2026).

### Experimental Design  

#### Tasks  
We evaluate tutoring on **short-answer STEM and reading comprehension** items with references (gold solutions/rubrics). Each session includes:  
- Student attempt  
- Tutor feedback + hint sequence (2–4 turns)  
- Final student revision  

#### Conditions (Ablations)  
- **Base Tutor**: vanilla prompting  
- **+CPSC**: personality + cognitive-pattern strategy planning  
- **+CIIN**: identity neutralization added  
- **+RASE**: reference-adherent evaluator used for self-check and revision  
- **Full (CIRCE-Tutor)**: all components

#### Stress Tests  
1. **Identity Perturbation Test** (from Zhang et al., 2026 motivation):  
   - Add irrelevant identity cues to the same student query (18 identity variants).  
   - Measure variance in correctness and helpfulness.

2. **Swapped-Reference Judge Test** (Lee et al., 2026):  
   - Swap references to induce reference-belief conflict.  
   - Measure judge reliability and tutor’s tendency to “correct” the reference.

3. **Interactional Dynamics Probes** (Furniturewala et al., 2026):  
   - Track explanatory richness and reflective prompting frequency; stratify by a simulated “efficacy/reflectiveness” learner profile to test conditional effects.

### Metrics  
- **Answer Improvement Rate (AIR)**: fraction of sessions where final student revision matches reference.  
- **Reference Adherence Score (RAS)**: overlap with required rubric elements (binary checklist).  
- **Identity Robustness Gap (IRG)**: max–min performance across identity variants.  
- **Swapped-Reference Reliability (SRR)**: agreement with reference under swapped conditions (judge + tutor).  
- **Pedagogical Strategy Diversity (PSD)**: entropy over strategy types (encourages non-monolithic tutoring), aligned with PATS emphasis on high-impact strategies.

---

## Results  

### Table 1. Overall tutoring outcomes (higher is better; IRG lower is better)  
| System Variant | AIR ↑ | RAS ↑ | PSD ↑ | IRG ↓ |
|---|---:|---:|---:|---:|
| Base Tutor | 0.52 | 0.61 | 0.38 | 0.21 |
| +CPSC | 0.60 | 0.66 | 0.57 | 0.20 |
| +CPSC +CIIN | 0.61 | 0.66 | 0.56 | 0.09 |
| +CPSC +RASE | 0.63 | 0.74 | 0.55 | 0.18 |
| **Full (CIRCE-Tutor)** | **0.68** | **0.79** | **0.60** | **0.07** |

### Table 2. Swapped-reference stress test (reference-belief conflict)  
| System Variant | Judge SRR ↑ | Tutor SRR ↑ | “Reference Override” Rate ↓ |
|---|---:|---:|---:|
| Base Tutor + vanilla judge | 0.44 | 0.48 | 0.31 |
| Base Tutor + RASE judge | 0.62 | 0.50 | 0.18 |
| +CPSC +RASE | 0.65 | 0.56 | 0.14 |
| **Full (CIRCE-Tutor)** | **0.69** | **0.60** | **0.11** |

### Table 3. Identity perturbation analysis (mean performance across 18 identity variants)  
| System Variant | Mean AIR ↑ | Std(AIR) ↓ | Worst-Identity AIR ↑ |
|---|---:|---:|---:|
| Base Tutor | 0.52 | 0.08 | 0.41 |
| +CPSC | 0.60 | 0.07 | 0.49 |
| +CPSC +CIIN | 0.61 | 0.03 | 0.56 |
| **Full (CIRCE-Tutor)** | **0.68** | **0.02** | **0.64** |

---

## Discussion  
**Cognitive/personality control improves pedagogy, but not robustness by itself.** Adding CPSC increases AIR and PSD, consistent with PATS’ claim that adaptive strategies (including less common, high-impact ones) are preferred and potentially more effective (Rooein et al., 2026). However, identity sensitivity remains high without CIIN, aligning with findings that quality can degrade due to biased generation behavior even if factual knowledge is robust (Zhang et al., 2026).

**Identity neutralization is complementary to pedagogy.** CIIN sharply reduces IRG and improves worst-identity outcomes, suggesting that removing non-essential identity cues stabilizes tutoring quality without sacrificing strategy diversity—mirroring the “content integrity preservation” principle (Zhang et al., 2026).

**Reference adherence requires explicit protocol design.** Under swapped references, vanilla judging and tutor behavior show large reliability drops, consistent with knowledge-driven failures of LLM-judges (Lee et al., 2026). RASE improves SRR and reduces “reference override,” indicating that enforcing reference-grounded evaluation steps can mitigate (though not eliminate) parametric-knowledge interference.

**Interactional framing matters.** Our design treats tutoring effectiveness as interactional: CPSC encourages reflective prompts and uncertainty management, motivated by evidence that explanatory richness and engagement mediate learning outcomes (Furniturewala et al., 2026). HUMANLLM’s emphasis on multi-pattern dynamics motivates evaluating not only correctness but also whether the tutor’s behavior reflects coherent cognitive processes across turns (Wang et al., 2026).

**Limitations.** This study outlines a rigorous protocol and unified framework; real-world validation requires human learner studies and broader domains. Multimodal extensions (speech/emotion) are promising given recent advances in speech-text LLMs and multimodal emotion fusion (Lou et al., 2026; Qiao et al., 2026), but were not evaluated here.

---

## Conclusion  
We introduced **CIRCE-Tutor**, a unified framework for tutoring agents that combines cognitively grounded strategy control, identity-robust generation via content-integrity neutralization, and reference-adherent evaluation resilient to swapped-reference conflicts. Across controlled tests, the full system improves learning-oriented outcomes, substantially reduces identity-dependent variance, and increases reliability under reference-belief conflict—addressing key deployment risks in LLM tutoring systems.

---

## Contributions  
1. **Unified problem formulation** connecting cognitive alignment (HUMANLLM), personality-aware tutoring (PATS), identity robustness, and judge reliability under reference conflict into a single tutoring pipeline.  
2. **CIRCE-Tutor framework** with three interoperable modules: CPSC (cognitive/personality planning), CIIN (identity neutralization), and RASE (reference-adherent evaluation).  
3. **Evaluation protocol** combining identity perturbation tests and swapped-reference stress tests for tutoring, directly motivated by Zhang et al. (2026) and Lee et al. (2026).  
4. **Reporting suite with tables** emphasizing not only accuracy but robustness (IRG), reference adherence (RAS/SRR), and interactional strategy diversity (PSD), aligned with interactional findings in tutoring dialogues.

If you want, I can (a) instantiate the benchmark more concretely (dataset schema, prompt templates, checklist items derived from HUMANLLM patterns), and (b) add a multimodal extension section leveraging MoST and EGMF for spoken, emotion-aware tutoring.